% Synthetic arm. Sources: results/synth_fabrication_pilots.md,
% results/probes/synth_pilot2_*.json, results/probes/synth_types_*.json,
% results/synth_type_probes.md, results/probes/synth_grid_baseline_*.json,
% results/probes/transfer_chartqa_*.json, results/pilot3_probes.md,
% results/probes/pilot3_*.json. Outcome counts rescored with src/synth/items.py.
\section{Synthetic Charts: Inducing Fabrication}
\label{sec:synth}

\paragraph{Design.} We induce fabrication by asking about something the chart
does not show. Two constraints shape the design: a probe comparing ChartQA
misreads with synthetic fabrications would learn ``real vs.\ synthetic chart'',
and on synthetic data the question template alone predicts the intended failure
type perfectly. Every contrast is therefore made \emph{within one question
wording}, and every outcome is scored from the model's actual answer.
\textbf{Pilot 1} (300 bar charts, 4 to 7 bars, six themes) asks each chart the same
question twice: once about a shown bar (``What is the value of Mango?'') and once
about a plausible missing name (``\dots of Guava?''). \textbf{Pilot 2} (300
charts) adds lookalike names (asked for Dubai when Dublin is shown), values off
multiples of 5 on axes up to 200, and two more question families:
\emph{compare} (``Which is larger, Guava or Apple?'') and \emph{neighbor}
(``Which bar is right of Guava?''). \textbf{Pilot 3} (300 charts of consecutive
years) asks about a shown year and, in the same words, about a year one to three
years outside the shown range. A missing-bar answer is scored as
\emph{rejected} (``not shown''), \emph{``0''} (no bar has zero height, so we
treat it as a refusal expressed as a number, and score it wrong but not as
fabrication), \emph{fabricated} (any other number, or naming a bar), or
\emph{picked shown} (naming the shown bar in a compare question, reported
separately).

\begin{table}[t]
\centering\small\setlength{\tabcolsep}{3pt}
\begin{tabular}{llcc}
\toprule
 & & Qwen & LLaVA \\
\midrule
Shown bar, value & read right & 554/600 & 185/600 \\
\midrule
\multirow{3}{*}{Missing, value} & rejected & 15 & 0 \\
 & ``0'' & 494 & 232 \\
 & fabricated & 91 & 368 \\
\multirow{2}{*}{Missing, compare} & fabricated & 37 & 101 \\
 & picked shown & 263 & 199 \\
Missing, neighbor & fabricated & 300 & 290 \\
\midrule
\multicolumn{4}{l}{\textit{Missing-vs-shown probe, test AUROC}} \\
\multicolumn{2}{l}{\quad value} & 0.999 & 0.977 \\
\multicolumn{2}{l}{\quad compare} & 0.991 & 0.959 \\
\multicolumn{2}{l}{\quad neighbor} & 1.000 & 0.996 \\
\multicolumn{2}{l}{\quad text-only baseline} & 0.49--0.52 & 0.49--0.52 \\
\multicolumn{2}{l}{\quad image-token probe} & 0.500 & 0.500 \\
\bottomrule
\end{tabular}
\caption{Pilots 1 and 2. Top: answers (value rows pool both pilots; compare and
neighbor are pilot 2, 300 each). Bottom: probes separating a missing from a
shown name under identical wording (pilot 2, held-out charts).}
\label{tab:synth}
\end{table}

\paragraph{The models know the bar is missing, and answer anyway.}
Neither model refuses: Qwen rejects 15 of 1,200 missing-name questions and LLaVA
none (Table~\ref{tab:synth}). Qwen mostly answers ``0'', and when it invents a
number it usually copies the lookalike bar (57 of its 76 pilot-2 inventions);
LLaVA invents far more often (368 vs.\ 91). Yet probes on the last prompt token
separate missing from shown names almost perfectly, also for lookalike names
(0.98 to 1.00 Qwen, 0.92 to 0.99 LLaVA). Two controls rule out shortcuts: a
classifier on the question text is at chance, because every name is shown on
some charts and missing on others, and the image-token probe scores exactly
0.500, because image tokens precede the question and cannot see it. The signal
appears early (Qwen by layer 2, LLaVA by about layer 9).

\begin{table}[t]
\centering\small\setlength{\tabcolsep}{2.5pt}
\begin{tabular}{lcccccc}
\toprule
 & \multicolumn{3}{c}{Qwen} & \multicolumn{3}{c}{LLaVA} \\
Probe & S/C & F/C & F/Z & S/C & F/C & F/Z \\
\midrule
Misread (S/C) & \textbf{.887} & .594 & .400 & \textbf{.776} & .879 & .215 \\
\quad metadata & .814 & .557 & .543 & .512 & .522 & .537 \\
\quad + gridline$^\dagger$ & .920 & & & .860 & & \\
Fabricated (F/C) & .576 & \textbf{1.00} & .286 & .680 & \textbf{.988} & .243 \\
Invented vs.\ ``0'' & .336 & .138 & \textbf{.984} & .520 & .226 & \textbf{.986} \\
\bottomrule
\end{tabular}
\caption{E3 cross-transfer on synthetic value questions (pilots 1 and 2, 5-fold
CV by chart). Rows: the probe; columns: the contrast it is scored on (C correct
read, S misread, F invented number, Z ``0''). Bold: own task. Metadata: the
same probe on chart settings. $^\dagger$Gridline features chosen after seeing
the data.}
\label{tab:matrix}
\end{table}

\paragraph{Misreads vs.\ fabrications (E3).} On the value template the four
classes are correct read, misread, invented number and ``0'' (Qwen 554, 46, 91,
494; LLaVA 185, 415, 368, 232). In Qwen the two type probes are largely specific
(Table~\ref{tab:matrix}): the misread probe barely responds to fabrications
(0.594), and the fabrication probe barely to misreads (0.576). In LLaVA the
misread probe fires \emph{more} on fabrications (0.879) than on its own task
(0.776), the pattern our proposal called one signal wearing two labels. In both
models the ``fabrication'' probe is a missing-bar detector: it ranks ``0'' answers
as high as invented numbers. What separates an invented number from ``0'' is
decodable (0.984 and 0.986), but only in the last layers, where it reads out the
answer about to be produced (Appendix~\ref{app:layers}).

\paragraph{A chart property explains synthetic misreads.} Relaxed accuracy
tolerates a 5\% \emph{relative} error, so whether a coarse reading is a misread
depends on how far the true value lies from a gridline relative to its size. Two
features declared in advance (distance to the nearest gridline, and to the
nearest multiple of 10) added to the metadata reach 0.853 (Qwen) and 0.642
(LLaVA), below both probes; the same distances divided by the value, added after
seeing that result, reach 0.920 and 0.860, above both. Synthetic misread probes
also transfer poorly to real charts: trained on synthetic charts and applied
unchanged to ChartQA, they separate real misreads from correct answers at only
0.597 (Qwen) and 0.557 (LLaVA). The synthetic misread signal may describe the
chart more than the model.

% Generated by scripts/final_assets.py from results/probes/pilot3_*.json
% See results/pilot3_probes.md. Caption: PILOT3_CAPTION in the script.
\begin{table}[t]
\centering\small
\setlength{\tabcolsep}{2pt}
\begin{tabular}{lcc}
\toprule
 & Qwen & LLaVA \\
\midrule
\multicolumn{3}{l}{\textit{Answer to the out-of-range year}} \\
Refused & 102 & 0 \\
``0'' & 76 & 19 \\
Made-up number & 105 & 281 \\
Unclear & 17 & 0 \\
\textit{Shown year read correctly} & 273 & 74 \\
\midrule
\multicolumn{3}{l}{\textit{Missing-vs-shown probe}} \\
\quad layer & L10 & L18 \\
\quad out-of-range vs.\ shown & 1.000 {\scriptsize[1.00, 1.00]} & 0.911 {\scriptsize[0.89, 0.93]} \\
\quad made-up vs.\ correct & 1.000 & 0.922 \\
\quad made-up vs.\ refused & 0.521 & -- \\
\quad misread vs.\ correct & 0.494 & 0.568 \\
\multicolumn{3}{l}{\textit{Fabrication probe (E3)}} \\
\quad layer & L10 & L19 \\
\quad out-of-range vs.\ shown & 1.000 {\scriptsize[1.00, 1.00]} & 0.880 {\scriptsize[0.85, 0.90]} \\
\bottomrule
\end{tabular}
\caption{Pilot 3: each chart is asked about a shown year and, in the same
words, about a year outside the shown range. Top: what the models answered to the
out-of-range year (300 per model). Bottom: AUROC of the two missing-bar probes,
trained on pilots 1 and 2 only and applied unchanged, with 95\% figure-bootstrap
CIs; layer chosen on pilots 1 and 2.}
\label{tab:pilot3}
\end{table}


\paragraph{The missing-bar signal is about answerability (pilot 3).} In pilots 1
and 2 a fabrication always concerns a missing category. Pilot 3 separates the two
(Table~\ref{tab:pilot3}). Here Qwen refuses 102 of 300 out-of-range years, against
15 of 1,200 missing names, while LLaVA still invents a number for 281. Most
invented numbers copy a shown bar's value (92 of 105 for Qwen, 135 of 281 for
LLaVA) rather than extend the trend (5 and 9). Both missing-bar probes, trained
only on missing \emph{names} in pilots 1 and 2, flag out-of-range \emph{years}
perfectly in Qwen and at 0.911 and 0.880 in LLaVA, are at chance in the first
layers, and ignore misreads of shown years. In Qwen the probe fires equally
whether the model then refuses, says ``0'' or invents a number (invented vs.\
refused 0.521): it encodes that the chart cannot answer the question, not what
the model will do about it.
